\section{新型装配模式与系统实现}
\subsection{新型装配模式设计}
在现有动力电池包装配系统中，关键瓶颈并不单纯来自设备能力，而在于知识与数据分散于不同环节。工艺规则掌握在工艺工程师手中，检测经验沉淀于质量人员经验中，设备能力保存在 PLC、机器人程序或视觉模型中，而现场异常往往依赖人工临时判断。其结果是，即使产线已经部署了机器人和视觉系统，在面对产品换型、异常返修或跨工位质量追溯时，仍然缺少统一的认知支撑。

基于这一问题，本文提出知识驱动的人机协同动力电池包智能装配与在线检测模式。该模式的基本原理是：将动力电池包的产品结构、零部件关系、工艺步骤、质量标准、设备能力、工位约束和异常案例以知识图谱的方式显式表达，并在其上叠加工艺规则推理与大模型智能体，实现面向装配现场的知识调用、异常解释和任务辅助决策 \cite{Sahota2025JIM,Meyers2024JIM,Wang2025CJME,Zhang2023JMSYKG,Chen2022AEI,He2023Access}。与传统以设备控制逻辑为中心的装配组织方式相比，该模式将“知识”从隐性经验转化为可调用、可追溯、可约束的系统资源，从而使装配过程不再只是按预设程序执行，而是能够围绕质量状态和工艺上下文进行动态调整。

与传统模式相比，这一新模式并不将大模型智能体置于执行层，而是将其放在认知辅助层。执行层仍由机器人、协作机器人、视觉检测站、焊接设备、力矩设备和人工工位构成；知识层由工艺知识图谱、质量规则库和设备能力模型构成；决策辅助层则由大模型智能体承担自然语言交互、工艺问答、异常定位建议和多工位协调建议。这一分层设计主要基于工业场景对稳定性、可解释性和安全性的要求：大模型智能体直接控制设备的风险较高，而将其用于知识检索、工艺解释和辅助决策更为现实，也更容易落地 \cite{Pulikottil2023IJAMT,Ji2024SciRep,Sgarbossa2024JMSY,Long2025IJPR}。因此，该模式不是用大模型智能体替代产线，而是在知识图谱约束下增强现有产线的柔性、可解释性和闭环能力。

图~\ref{fig:newmode-hrc-workflow}给出了嵌入算法后的人机协同系统工作流。图中从语音输入、环境描述、视觉基础模型感知到大模型任务推理、代码生成和机器人执行形成了一条连续链路。尽管原图并非针对动力电池包场景，但它所体现的“多模态输入-算法推理-协同执行”组织方式，与本文提出的新模式具有较强对应性，也更具体地支撑了前文关于认知辅助层和执行层分工的讨论 \cite{Ji2024SciRep}。\vspace{-0.4em}

\begin{figure}[H]
  \centering
  \includegraphics[width=0.62\textwidth]{fig8_hrc_workflow.png}
  \caption{语言驱动的人机协同工作流}
  \label{fig:newmode-hrc-workflow}
\end{figure}

从功能上看，该模式至少包含四项改进。第一，工艺知识不再以文档和口头经验分散存在，而是通过知识图谱统一表达，便于搜索、复用和约束推理。第二，在线检测结果不再只是合格/不合格标签，而是成为驱动异常诊断和后续工艺调整的数据输入。第三，协作工位中的操作员可以通过自然语言与大模型智能体交互，快速获取“当前故障可能由哪些工序引起”“返修顺序如何安排”“该型号与上一型号的工艺差异是什么”等答案。第四，系统可以将换型、返修和多工位资源分配问题交由多智能体规划框架辅助处理，从而减少人工调度负担 \cite{Massouh2025JIM,Abelha2020RAL,Yao2023JMSY}。

这一模式的价值不在于单独引入知识图谱或大模型智能体，而在于将它们与动力电池包的具体工艺相结合。其核心思路是由设备负责执行，视觉与传感器负责感知，知识图谱负责组织工艺与约束，大模型智能体负责解释、检索和辅助决策，人则保留对关键风险和异常处置的最终把关。这样，模式设计、基本原理、系统实现和应用评价之间便形成了较清晰的对应关系。
